A l'abadia de Westminster, a Anglaterra, s'hi coronen les reines i els reis britànics i s'hi enterren o honoren els personatges més distingits. També s'hi ret homenatge a una equació, l'equació de Dirac, que data de l'any 1928 i és l'única que hi ha per ara. L'equació de Dirac és un exemple d'equació amb conseqüències inesperades. Amb aquesta equació es va predir l'antimatèria a partir de la deducció de l'existència de l'antielectró, que després va passar a anomenar-se \textit{positró}. El positró va ser descobert experimentalment el 1932, un any després d'haver estat predit.

\begin{apartat}{1,25}
Quina és l'energia mínima necessària d'un fotó, en MeV, perquè es materialitzi en un parell electró-positró? Calculeu la freqüència, la longitud d'ona i la quantitat de moviment d'aquest fotó.
\end{apartat}

\begin{apartat}{1,25}
En la desintegració $\beta^+$, es crea una partícula que és antimatèria. De quina partícula es tracta? En la desintegració $\beta^-$, quina partícula que també és antimatèria es crea? En el procés d'anihilació positró-electró, per quin motiu s'han de crear un parell de fotons que tenen la mateixa energia que el sistema i que viatgen en sentits oposats? Quant val la càrrega inicial, abans de l'anihilació positró-electró? I la càrrega final, després de l'anihilació?
\end{apartat}

\dades{%
  Massa de l'electró = massa del positró $= 9,11\times 10^{-31}\ \mathrm{kg}$.\\
  $1\ \mathrm{eV} = 1,602\times 10^{-19}\ \mathrm{J}$.\\
  $c = 3,00\times 10^{8}\ \mathrm{m\,s^{-1}}$.\\
  $h = 6,63\times 10^{-34}\ \mathrm{J\,s}$.}
